\subsection{Strange Modules}\label{strange-modules}


\subsection{Questions (and some answers)}
\begin{enumerate}
\item Is dualization an anti-equivalence of finitely adequate modules?

      Yes.
\item Are finitely adequate modules closed under extension?

      If the quotient module is also finitely generated, then yes.
\item\label{ad-hom-closed} Are finitely adequate modules closed under $\Hom$?

      Yes, \Cref{prp:fin-adeq-hom-closure}.
\item Are finitely adequate modules closed under $\_\otimes\_$?

      Yes, for finite free \(B\), by \Cref{ad-hom-closed} and \(A\otimes B^\ast=\Hom(B,A)\), but no in general:


\item Does exponentiating with closed propositions preserve finitely adequate?

      Yes, by partial answer to \Cref{finite-exp-closed}.
\item\label{finite-exp-closed} Are projective/proper/finite-scheme products of finitely adequate modules finitely adequate?

      Partial answer:
      \begin{enumerate}[(i)]
      \item Let $X$ be a finite scheme, i.e. $X$ an affine scheme such that $R^X$ is a finitely presented $R$-module.
      For any finitely adequate-module $M:X\to\Mod{R}_{\mathrm{ad}}$ on $X$, $(x:X)\to M_x$ is finitely adequate: We can exponentiate the diagram witnessing that it is an finitely adequate-module and we can globalize the presentations of pointwise fp modules.
      Exponentiating is left exact and since $M$ is wqc and $X$ affine, it is also right exact.
      \end{enumerate}

\item Are the cohomology groups of finitely adequate module bundles on projective schemes finitely adequate?

      None so far. It seems a bit much to ask, but it would be really great for Serre-Duality, so it should be worth it to try to adapt e.g.\ \cite[19.1.3]{vakil}.

\item Is any finitely adequate-module quasi-coherent?

      No. The kernel $K$ of $2\cdot\_:R\to R$ can not be (strongly) quasi-coherent in general.
      So far there is only an external argument: \url{https://www.youtube.com/watch?v=mdjq91Gc6Lo} around 1:10:00 states that strongly quasi-coherent modules are invariant under pullback. But in general $K$ can change between trivial and non-trivial on different stages.
\item Does being finitely adequate depend only on the underlying type?  Maybe related to being an fppf sheaf?
\item What are examples of modules that are not finitely adequate?

      $R^\mathbb{N}$.  A direct sum of countable copies of $R$.  (Citations/proofs?)

      Are there some that are ``smaller'' in some sense?
      What about the ones mentioned in \Cref{strange-modules}?


\item What are examples of modules that are not infinitesimally linear/tangential?

      $\mathrm{Nil}(R)$ is infinitesimally linear but not tangential.

      Denote by $R[S]$ the free $R$-module generated by a set $S$.
      Let $M :\equiv R[\D(2)]/R[\D(1)\vee\D(1)]$ and say $\eta : \D(2) \to R[\D(2)]$ is the constructor map, which induces a pointed map $\tilde{\eta} : \D(2) \to M$.
      By construction we have $\tilde{\eta}|_{\{0\} \times \D(1)} = \tilde{\eta}|_{\D(1) \times \{0\}} = 0$ but since $\D(1)\vee\D(1) \lhook\joinrel\to \D(2)$ is provably not surjective, we have $\tilde{\eta} \neq 0$.
      Hence $M$ is not infinitesimally linear (and in particular not finitely adequate).

      If $R$ has characteristic $p$, then $\hood{R,0}{p-1}$ is not tangential. Its tangent space is $R$ by the same above reasoning.

\item What is the relation of finite adequacy and finitely presentedness/linear adequacy?

      \begin{lemma}
        Not every finitely adequate $R$-module is either finitely presented or linearly adequate.
        Specifically, this is not the case for the family of ideals $(r)$, for $r:R$.
      \end{lemma}

      \begin{proof}
        We first show that if $(r)$ is linearly adequate, then $(r=0) \lor (r \neq 0)$. Indeed, linearly adequate modules are injective in the category of finitely adequate $R$-modules, so we get a lift
        \[
        \begin{tikzcd}
        & R \ar[d, "\cdot r", twoheadrightarrow] \\
        (r) \ar[r, "id"]\ar[ur, "\exists s", dashed] & (r)
        \end{tikzcd}
        \]
        meaning that $(r)$ is a direct summand of $R$ such that $(r) = (e)$ for an idempotent $e:R$.
        Since $R$ is local, $(e=0) \lor (e=1)$ and thus $(r=0) \lor (r \neq 0)$.

        A dual argument shows that if (r) is finitely presented, then also $(r=0) \lor (r \neq 0)$.
        But assuming that all $(r)$ were to admit either one of these presentations, we would get $\prod_{r:R} (r=0) \lor (r \neq 0)$, which is false.
      \end{proof}

      \begin{remark}
        This lemma cannot be improved in the sense that here is no finitely adequate $R$-module $M$ that is not finitely presented.
        Indeed, if $M$ were not finitely presented, then it would be not finite free.
        But finitely adequate modules are \emph{not not} finite free.
        A similar statement holds for linearly adequate.
      \end{remark}

      \item
      \begin{remark}\label{rem:hood-module-prs-surj}
        If $f : M \rightarrow N$ be an epimorphism of $R$-modules with $M$ linearly adequate
        and $N$ finitely adequate,
        it does not follow that $\hood{f,0}{1}$ is surjective.
        Here is an example.

        Let $\varepsilon : \D(1)$ be arbitrary.
        Consider the map $\pi_\varepsilon : R \to (\varepsilon)$ which sends $r$ to $r \varepsilon$.
        Note that $R$ is linearly adequate and $(\varepsilon)$ is finitely adequate
        by \Cref{thm:fin-adeq-class}.
        We will show that $\pi_\varepsilon$ does not satisfy the claim for all $\varepsilon$.
        Indeed, assume that $\hood{\pi_{\varepsilon}, 0}{1}$ is surjective for all $\varepsilon$.
        Note that $\varepsilon \in \hood{(\varepsilon), 0}{1} = (\varepsilon)$.
        Let $r:\D(1) = \hood{R,0}{1}$ be a lift of $\varepsilon$.
        Since $\pi_\varepsilon(1) = \varepsilon$, $\pi_\varepsilon(r-1)=0$, i.e., $(r-1)\varepsilon=0$.
        But $r$ is nilpotent, hence $r-1$ invertible and this would mean $\varepsilon = 0$, which cannot be true for all $\varepsilon$.
      \end{remark}
\end{enumerate}
